% Generated draft: requires mathematical and source review.
\documentclass[11pt]{article}
\usepackage[margin=25mm]{geometry}
\usepackage{fontspec}
\setmainfont{Latin Modern Roman}
\usepackage{amsmath,amssymb,mathtools,unicode-math}
\setmathfont{Latin Modern Math}
\usepackage{enumitem,needspace,xurl,hyperref}
\setlength{\parindent}{0pt}
\setlength{\parskip}{5pt}
\setlength{\emergencystretch}{3em}
\newcommand{\N}{\mathbb{N}}
\newcommand{\Z}{\mathbb{Z}}
\newcommand{\Q}{\mathbb{Q}}
\newcommand{\R}{\mathbb{R}}
\newcommand{\C}{\mathbb{C}}
\newcommand{\F}{\mathbb{F}}
\newcommand{\A}{\mathbb{A}}
\newcommand{\PP}{\mathbb{P}}
\newcommand{\E}{\mathbb{E}}
\newcommand{\eps}{\varepsilon}
\newcommand{\dd}{\,\mathrm{d}}
\begin{document}
\section*{UnsolvedMath ID: 10400180}
\textbf{Source problem code:} catalogue entry
\subsection{Definitions and mathematical statement}
For an oriented closed \(3\)-manifold \(M\), a \(\operatorname{Spin}^c\) structure \(\alpha\) is a lift of its oriented tangent-frame bundle to
\[
 \operatorname{Spin}^c(3)
 =\bigl(\operatorname{Spin}(3)\times U(1)\bigr)/\{(1,1),(-1,-1)\}.
\]
Pairs are considered up to orientation-preserving diffeomorphisms carrying one \(\operatorname{Spin}^c\) structure to the other, in the conventional oriented interpretation. The source itself says only ``closed \(3\)-manifold'' and leaves orientation conventions implicit.

\emph{Question (M. Polyak).} Is there a ``Rokhlin invariant'' of a pair \((M,\alpha)\), where \(M\) is a closed \(3\)-manifold and \(\alpha\) is a \(\operatorname{Spin}^c\) structure on \(M\)? See the source's Question 10.21, whose statement is not supplied.

For comparison, in the conventional spin setting the Rokhlin invariant of an oriented spin \(3\)-manifold \((M,s)\) is described by \(\sigma(W)\bmod16\) for a smooth compact spin \(4\)-manifold extending \((M,s)\), with the usual normalization. The quotation marks in the source do not specify which feature of this construction is to be generalized. A proposed \(\operatorname{Spin}^c\) analogue must give its target, normalization, defining data, independence of auxiliary choices, and intended relation to the spin invariant. The excerpt supplies none of these requirements in precise form. In particular, it does not explicitly demand that forgetting a spin structure to its associated \(\operatorname{Spin}^c\) structure preserve every spin Rokhlin value; such a compatibility condition is not silently added.
\subsection{Short English statement}
Polyak asks for an analogue of the Rokhlin invariant attached to a closed \(3\)-manifold with a \(\operatorname{Spin}^c\) structure. The value group, normalization, and required relationship with the spin Rokhlin invariant are unspecified, and the cited Question 10.21 is absent.
\subsection{Sources}
\begin{itemize}
\item Local UnsolvedMath registry, permanent ID 10400180, source code catalogue entry; attribution M. Polyak, Question 10.5. The supplied title is truncated after ``(M.''.
\item Ohtsuki, \emph{Problems on invariants of knots and 3-manifolds} (2002), Question 10.5, PDF page 145; explicit cross-reference to Question 10.21. Supplied URL: \url{https://msp.org/gtm/2002/04/gtm-2002-04-024s.pdf}; recorded access 2026-07-29, extraction pdf-font-structure. The collection distinguishes 231 numbered bold problem, conjecture, and question headings.
\item Supplied status open and OPEN-TRIAGE. The research summary identifies Polyak's question about \(\operatorname{Spin}^c\) structures. Current primary-literature status and redistribution terms for the publicly accessible author/publisher copy remain NEEDS\_REVIEW.
\end{itemize}
\end{document}
